\section{Scope, interpretation, and implications}

\subsection{What the dynamic smoothness object means}

The selected smoothness is not treated as a fixed property of a series. At a given origin, several local minima of the forecast-loss surface may coexist. Their temporal branches summarize recurring smoothing regimes discovered by the data.

For a final selected value $\widehat S_T$, the effective-degrees-of-freedom interpretation remains
\begin{equation}
\label{eq:scope_smoothness_edf}
\widehat S_T
=
\frac{L-\operatorname{edf}_T}{L-d}.
\end{equation}
A value near zero corresponds to a flexible trend; a value near one corresponds to a trend close to the null-space polynomial. The dynamic contribution is that today's value can depend on the history of a persistent local-minimum branch rather than on one pooled optimum alone.

\subsection{Branches are data-driven, not fixed smoothness bins}

The continuation radius $\varepsilon$ does not define a preassigned category such as ``smoothness around 0.9.'' Branch locations arise from local minima of successive forecast-loss surfaces. The radius only determines whether two neighboring minima are close enough to be interpreted as the same branch.

Consequently, a branch can drift over time. Its identity reflects local continuity of the optimum, not membership in a fixed interval.

\subsection{What may be averaged}

The framework distinguishes three objects that should not be conflated. Historical forecast losses may be averaged to rank candidate smoothness values or branches. Historical smoothness values may be averaged when a declared branch functional such as a recent-mean rule is used. Historical fitted trends are never averaged in the proposed procedure.

After the current smoothness is selected, the trend is always refit on the newest information available before the forecast.

\subsection{Static and dynamic targets}

The pooled selector asks which single value of $S$ performs best on average over historical forecast origins. The branch formulation asks how a time-indexed collection of local minima can be organized and converted into multiple admissible current values of $S$. Predictive improvement is a possible outcome of particular decisions, not a property guaranteed by the representation itself.

Neither target is automatically the smoothness that best reconstructs a latent historical trend. Forecasting and signal recovery remain distinct objectives.

\subsection{Numerical boundary}

The forecasting paper does not claim that a particular root-search or assignment algorithm is the substantive contribution. Recovery of multiple minima, exact endpoint handling, and temporal correspondence of minima belong to the companion numerical-methods paper.


\subsection{Boundary behavior of smoothness forecasts}

Normalized smoothness is restricted to $[0,1]$. Convex means of historically
admissible values remain in that interval, whereas linear or increment
extrapolations of the branch trajectory need not. Truncation back to the
admissible interval is therefore part of defining these functionals,
not an innocuous plotting convention. The simulations show that the
frequency of truncation varies substantially with the form of $\phi$.
This is a concrete difference between equally admissible constructions
of the next smoothness value.

\subsection{Limits of predictive interpretation}

A favorable four-series confirmation result is not sufficient to establish
universal superiority. The broad financial panel and controlled changing-
roughness experiments did not reproduce an aggregate advantage over
pooled forecast-CV. Their errors are evaluated on different scales, and
the series within the financial panel are not independent.
Accordingly, the paper claims a reusable representation for chronologically
tracked forecast-loss minima and an illustrative family of decision maps;
it does not claim that tracked branches inevitably improve forecasting.

\subsection{Limitations}

The baseline estimator assumes equally spaced observations, a fixed difference order, a fixed final estimation window, and deterministic continuation of the estimated trend. Checkpoints 05 and 06 show that high-order native polynomial continuation can amplify small differences in selected $S$ into extreme forecast errors. In particular, restricting higher difference orders reduced the catastrophic error tail in a post-hoc diagnostic, but does not validate a revised forecasting selector. The current branch construction also introduces design quantities such as the tracking radius, recent-history length, and weighting parameters. Those quantities must be selected on development data and frozen before outer-test evaluation.

Local-minimum identity can become ambiguous when branches approach or cross. The initial greedy one-to-one matcher is therefore a transparent baseline rather than a mathematical definition of true branch identity. Sensitivity to the numerical correspondence rule must be reported separately from forecasting performance.

Finally, the framework isolates trend extrapolation. It does not yet add a separate stochastic residual forecast model. That restriction keeps the interpretation of any forecast gain tied to smoothness selection rather than to another predictive component.
